\section{预势与瞬子计算}\label{sec:instcount}
在第 \ref{sec:purecurve} 小节中，我们学会了如何借助 Seiberg-Witten 曲线求出纯 $\SU(2)$ 理论的 $a$ 和 $a_D$。将 $a_D$ 积分两次，便可以具体确定预势 $F(a)$。在本附录中，我们先进行这一计算，然后非常简要地解释，如何通过微观路径积分计算得到同一个预势。

\subsection{由曲线求预势}\index{Prepotential@预势|hyperpage}\index{RenormalizationGroupRelation@重整化群关系|hyperpage}\index{RiemannReciprocityLaw@Riemann 互反律|hyperpage}
纯 $\SU(2)$ 理论的曲线已在式 \eqref{curve_pure} 中给出，这里再次写出：
\begin{equation}
\Lambda^2 z+\frac{\Lambda^2}z=x^2-u,
\end{equation}
其微分为 $\lambda=xdz/z$。$a$ 和 $a_D$ 由下式确定：
\begin{equation}
a=\frac{1}{2\pi i}\oint_A\lambda,\qquad
a_D=\frac{1}{2\pi i}\oint_B\lambda.\label{aad}
\end{equation}
由此，应当确定满足 $a_D=\partial F/\partial a$ 的 $F(a)$。低能耦合为 $\tau(a)=\partial a_D/\partial a$。在式 \eqref{tauexpansion} 中，我们得到了 $\tau(a)$ 的弱耦合展开形式。积分两次，便看到 $F(a)$ 具有展开
\begin{equation}
2\pi iF(a)=-4a^2\log\frac{a}{\Lambda}+\sum_{k=1}^\infty d_k\frac{\Lambda^{4k}}{a^{4k-2}}\label{Fweak}
\end{equation}
其中 $d_k$ 是纯数。下面遵循文献 \cite{Chan:1999gj}，明确求出它们。

为此，使用所谓的重整化群关系
\begin{equation}
2\pi i\Lambda\frac{\partial}{\partial\Lambda}F(a,\Lambda)=4u\label{matone}
\end{equation}
它最初在文献 \cite{Matone:1995rx} 中推导出来。这一关系可以用多种方法证明。一种直接的方法是注意到，由式 \eqref{lambdaSU2pure} 可得
\begin{equation}
2\pi i\Lambda\frac{\partial}{\partial\Lambda}=4\frac{\partial}{\partial\tau_{UV}}
\end{equation}
这里保持紫外调节参数 $\Lambda_{UV}$ 不变。回顾紫外能标处的预势为 $F(\phi)=(\tau_{UV}/2)\tr\phi^2$，见第 \ref{sec:lel} 小节末尾。于是，关系 \eqref{matone} 只是下述陈述：
\begin{equation}
\frac{\partial}{\partial\tau_{UV}}\vev{F}=\frac12\vev{\tr\phi^2}=u.
\end{equation}

也可以选择从式 \eqref{aad} 出发，以纯数学方式证明关系 \eqref{matone}。从显然的关系开始：
\begin{equation}
(\Lambda\frac{\partial}{\partial\Lambda}+a\frac{\partial}{\partial a})F=2F.
\end{equation}
利用它，得到
\begin{align}
\frac{\partial}{\partial u}\Lambda\frac{\partial}{\partial\Lambda}F
&=\frac{\partial}{\partial u}(2F-a\frac{\partial}{\partial a}F)
=a\frac{\partial a_D}{\partial u}-\frac{\partial a}{\partial u}a_D\\
&=\frac1{(2\pi i)^2}\left(\oint_A\lambda\oint_B\frac{\partial}{\partial u}\lambda-\oint_A\lambda\oint_B\frac{\partial}{\partial u}\lambda\right).
\end{align}
这里，$\lambda=xdz/z$ 是第二类微分（即留数为零的亚纯微分），而 $\partial\lambda/\partial u=-dz/(2xz)$ 是第一类微分（即全纯微分）。因此，可以利用 Riemann 互反律来计算，例如见文献 \cite{GH} 第 2.3 节。具体地说，$\lambda$ 的二阶极点位于 $z=0$ 和 $\infty$。在这两点处，$w:=1/x$ 都是良好的局部坐标，并且有展开
\begin{equation}
\lambda=x\frac{dz}z=-2\frac{dw}{w^2}+\cdots,\qquad
\frac{\partial}{\partial u}\lambda=-\frac{dz}{2xz}=dw+\cdots.
\end{equation}
所以有
\begin{equation}
\frac{\partial}{\partial u}\Lambda\frac{\partial}{\partial\Lambda}F=\frac4{2\pi i},
\end{equation}
从而证明了式 \eqref{matone}。

现在回到式 \eqref{Fweak} 中系数 $d_k$ 的计算。令 $u=\underline a^2$，则
\begin{align}
a(\underline a)&=\frac1{2\pi i}\oint_A\lambda=\frac1{2\pi i}\oint_A(\sum_{l=0}^\infty
(z+\frac1z)^l\frac{(-1)^l(2l)!}{(1-2l)(l!)^2 2^{2l}}
\frac{\Lambda^{2l}}{\underline a^{2l-1}})\frac{dz}z\\
&=\underline a\sum_{k=0}^\infty\frac{\Lambda^{4k}}{\underline a^{4k}}\frac1{(1-4k)2^{4k}}\frac{(4k)!}{(2k)!k!k!}.\label{113}
\end{align}
现在，将展开 \eqref{Fweak} 代入关系 \eqref{matone}，得到
\begin{equation}
a^2(1+\sum_{k=1}^\infty k d_k\frac{\Lambda^{4k}}{a^{4k}})=\underline a^2.\label{114}
\end{equation}
关系 \eqref{113} 和 \eqref{114} 合在一起，足以递归确定 $d_k$。最前面的几项是
\begin{equation}
d_1=\frac12,\quad d_2=\frac5{64},\quad d_3=\frac3{64},\quad d_4=\frac{1469}{32768},\quad\ldots.\label{results-from-curve}
\end{equation}

\subsection{由瞬子计算求预势}
这里从更微观的观点，介绍同一预势的另一种计算。这一计算可以完全用纯 $\SU(2)$ 规范理论的路径积分来表述 \cite{Nekrasov:2002qd}。利用一个密切相关的五维场论，基本相同的计算可以用更物理的方式来表述。我们在这里采用后一种方法。这里的介绍将非常简略；更多细节请读者参阅文献 \cite{TachikawaInstantonReview}。

\subsubsection{$\Omega$ 背景}\index{OmegaBackground@$\Omega$ 背景|hyperpage}\index{NekrasovPartitionFunction@Nekrasov 配分函数|hyperpage}
首先利用这样一个事实：给定一个四维 $\cN{=}2$ 超对称规范理论，其规范群为 $G$，超多重态属于 $G$ 的表示 $R$，也可以考虑一个五维超对称规范理论，具有相同的规范群 $G$，超多重态属于相同的表示 $R$；将该理论紧致化在半径为 $\beta$ 的 $S^1$ 上，再取零半径极限 $\beta\to0$，就恢复原来的四维理论。

\begin{figure}[h]
\[
\inc{omega}
\]
\caption{五维 $\Omega$ 背景的示意图。\label{fig:omega}}
\end{figure}
我们在所谓的 $\Omega$ 背景上考虑这个五维理论，该背景按如下方式得到。首先，将 $\bR^5$ 看作 $\bC^2\times\bR$，坐标为 $(z_1,z_2,x_5)$，其中 $z_{1,2}$ 为复数，$x_5$ 为实数。然后作等同
\begin{equation}
(z_1,z_2,x_5)\sim(z_1e^{i\beta\epsilon_1},z_2e^{i\beta\epsilon_2},x_5+\beta).\label{omegaident}
\end{equation}
为了保持超对称性，必须在这一等同 \eqref{omegaident} 的同时进行 $\SU(2)_R$ 旋转
\begin{equation}
\diag(e^{i\beta(\epsilon_1+\epsilon_2)/2},e^{-i\beta(\epsilon_1+\epsilon_2)/2})\in\SU(2)_R.
\end{equation}
在作等同 \eqref{omegaident} 时，也可以选择进行规范旋转
\begin{equation}
e^{i\beta\vec a}\in G
\end{equation}
其中 $\vec a$ 是 $G$ 的李代数中的一个元素。图 \ref{fig:omega} 给出了该系统的示意图。

这个时空的第五方向为 $\beta$，并伴随着 $\bR^4$ 一侧转过角度 $\beta\epsilon_{1,2}$。当把 $x_5$ 看作时间方向时，系统被迫旋转，远离 $\bC^2$ 原点的任何激发都要耗费很大的能量。这使系统在有效意义上具有有限体积，大小为 $\sim1/(\epsilon_1\epsilon_2)$。因此，系统具有有限的配分函数 $Z(\epsilon_{1,2},\tau_{UV};\vec a)$，其对数表现为
\begin{equation}
\log Z(\epsilon_{1,2},\tau_{UV};\vec a)=\frac1{\epsilon_1\epsilon_2}\cF(\tau_{UV};\vec a)+\text{关于 $\epsilon_{1,2}$ 奇异性较弱的项},
\end{equation}
其中 $\cF$ 是系统单位体积的自由能。这里用五维耦合常数 $g_{UV5}$ 定义 $\tau_{UV}$；$g_{UV5}$ 的定义使规范动能项具有形式
\begin{equation}
\int dx_5\int d^4x\frac1{2g_{UV5}{}^2}\tr F_{\mu\nu}F_{\mu\nu}.\label{5dgaugekin}
\end{equation}
先对 $x_5$ 积分，产生一个因子 $\beta$，于是得到
\begin{equation}
\tau_{UV}=\beta\frac{4\pi i}{g_{UV5}{}^2}.
\end{equation}

一个基本事实是，这个自由能 $\cF$ 等于预势：
\begin{equation}
2\pi iF(\tau_{UV};\vec a)=\cF(\tau_{UV};\vec a)=\lim_{\epsilon_{1,2}\to0}\epsilon_1\epsilon_2\log Z(\tau_{UV};\vec a).\label{main-instanton-formula}
\end{equation}
取一个受控的极限 $\beta\to0$，便可以提取纯四维理论的预势。这一适用于任意 $\cN{=}2$ 理论的普遍事实，可以用多种方法推导，见文献 \cite{Nekrasov:2002qd} 或文献 \cite{Tachikawa:2004ur} 的附录。

\subsubsection{约化为超对称量子力学}\index{InstantonModuliSpace@瞬子模空间|hyperpage}\index{SupersymmetricQuantumMechanics@超对称量子力学|hyperpage}
这种重新表述听起来或许繁琐，但其好处在于，预势的计算如今归结为 $\Omega$ 背景 \eqref{omegaident} 中系统的配分函数 $Z$ 的计算。将系统约化为超对称量子力学，便可以用与第 \ref{sec:box} 小节中分析盒子内的纯 $\cN{=}1$ 规范理论大致相同的方法完成计算。注意，这里将额外维度 $x^5$ 看作时间方向 $t$。

总配分函数 $Z$ 是微扰部分 $Z_\text{pert}$ 与非微扰部分 $Z_\text{non-pert}$ 的乘积。第一部分只再现式 \eqref{Fweak} 中领头的对数项，其再现方式本身颇为有趣，但这里不再介绍。相反，下面集中讨论非微扰部分 $Z_\text{non-pert}$。计算这一部分时，可以像第 \ref{sec:box} 小节那样，限制在每个拓扑扇区中的最低能态。从式 \eqref{5dgaugekin} 可以看出，一个规范场构型的能量包含贡献
\begin{equation}
\int d^4x\frac1{2g_{UV5}{}^2}\tr F_{\mu\nu}F_{\mu\nu},\label{5dprefactor}
\end{equation}
正如式 \eqref{bound} 所示，其下界为 $8\pi^2|k|/g_{UV5}{}^2$，其中 $k$ 为瞬子数。

因此，唯一对 $Z$ 的计算有贡献的经典场构型，是在每个等时面 $t=\text{常数}$ 上都具有瞬子构型的场构型。瞬子数 $k$ 是整数，所以沿时间演化时不能改变。然而，固定 $k$ 的瞬子构成一个连续族，因此瞬子的形状可以变化。用 $M_k$ 表示这一连续族的参数空间，称为瞬子模空间。给定一点 $p\in M_k$，就有一个瞬子解 $F_{\mu\nu}(x_{1,2,3,4};p)$，其中 $\mu,\nu=1,2,3,4$。随时间 $t=x_5$ 变化，瞬子的形状也发生变化。换言之，存在 $M_k$ 中的一条路径 $p(t)$。此时，五维作用量 \eqref{5dgaugekin} 近似变为
\begin{equation}
\int dt\int d^4x\frac1{2g_{UV5}{}^2}\tr F_{\mu\nu}F_{\mu\nu}
=\int dt(\frac{8\pi^2}{2g_{UV5}{}^2}+G_{IJ}(p)\partial_t p^I(t)\partial_t p^J(t))\label{sqm}
\end{equation}
其中，右侧的第一项和第二项分别来自左侧时空指标 $\mu,\nu$ 中不含 $5$ 和含有 $5$ 的项。这里，指标 $I=1,\ldots,\dim M_k$ 对应于 $M_k$ 的局部坐标，而 $G_{IJ}(p)$ 是 $M_k$ 上的度规。

对于 $G=\SU(2)$，已知模空间 $M_k$ 是一个 $4|k|$ 维空间。因此，给定 $k$ 的系统约化为一个在 $M_k$ 上运动的点粒子的超对称量子力学；其玻色部分的作用量由式 \eqref{sqm} 给出。这本质上是 Born-Oppenheimer 近似；一个重要事实是，在超对称背景下，这一近似实际上是精确的。

在式 \eqref{bound} 附近，我们解释了 $k\ge0$ 和 $k\le0$ 的瞬子保持不同的超荷。因此，只有 $k\ge0$ 的瞬子对 $Z$ 的计算有贡献，而 $k\le0$ 的瞬子对 $\bar Z$ 有贡献。由这些考虑，得到下式：
\begin{equation}
Z_\text{non-pert}=\sum_{k=0}^\infty Z_k.
\end{equation}
这里，$Z_k$ 是 $M_k$ 上超对称量子力学的配分函数，其玻色部分的作用量为式 \eqref{sqm}。由于式 \eqref{sqm} 的第一项只给出能量的常数平移，适宜将其分离出来，写成
\begin{equation}
Z_k=e^{2\pi i\tau_{UV}k}\tilde Z_k\label{constantshift}
\end{equation}
其中 $\tilde Z_k$ 是 $M_k$ 上最低能态能量为零的超对称量子力学的配分函数。

\subsubsection{具体计算}\index{InstantonCounting@瞬子计数|hyperpage}\index{YoungDiagram@杨图|hyperpage}\index{ArmLegLength@臂长与腿长|hyperpage}
下面计算 $Z_k$。从最简单的情形 $k=0$ 开始。在此情况下，模空间 $M_0$ 只是单独一个点，而在一个点中运动的点粒子实在无事可做。希尔伯特空间只有一维，哈密顿量为零，因此有
\begin{equation}
Z_0=1.
\end{equation}

接着考虑 $k=1$。在此情况下，瞬子模空间为 $\bC^2\times(\bC^2/\bZ_2)$。其中，第一个 $\bC^2$ 指定瞬子在 $\bC^2$ 上的质心位置。一个 $k=1$ 的 $\SU(2)$ 瞬子还由大小 $\rho$ 和群流形中的一个规范方向来指定，其中 $\rho$ 是非负实数。由于规范场取值于伴随表示，规范方向取值于 $\SO(3)$ 群流形，即 $S^3/\bZ_2$。它们与大小 $\rho$ 一起参数化 $\bC^2/\bZ_2$。

为 $\bC^2$ 引入坐标 $(z_1,z_2)$，为 $\bC^2/\bZ_2$ 引入坐标 $(u,v)$，并作等同 $(u,v)\simeq-(u,v)$。在角度为 $\beta\epsilon_{1,2}$ 的旋转下，它们变换为
\begin{equation}
\begin{aligned}
(z_1,z_2,u,v)\mapsto{}&(e^{i\beta\epsilon_1}z_1,\epsilon_{i\beta\epsilon_2}z_2,\\
&e^{i\beta(\epsilon_1+\epsilon_2)/2}u,e^{i\beta(\epsilon_1+\epsilon_2)/2}v).
\end{aligned}
% [Tutor Note: p.178 式 (A.2.13) 原文第二个分量印作 $\epsilon_{i\beta\epsilon_2}z_2$，按上下文（与第一个分量
%  $e^{i\beta\epsilon_1}z_1$ 对称）应为 $e^{i\beta\epsilon_2}z_2$；正文照原文排版。]
\end{equation}
对 $z_1$ 和 $z_2$ 的作用直接来自定义；对 $u$ 和 $v$ 的作用则较为微妙。在规范旋转
\begin{equation}
\diag(e^{i\beta a},e^{-i\beta a})\in\SU(2)
\end{equation}
下，$z_{1,2}$ 不变，而
\begin{equation}
(u,v)\mapsto(e^{i\beta a}u,e^{-i\beta a}v).
\end{equation}

这些数据足以计算 $Z_2$。
% [Tutor Note: p.178 原文此处写 "compute $Z_2$"，但本段讨论的是 $k=1$ 的单瞬子贡献（下一段才计算
%  $Z_2$），按上下文应为 $Z_1$；正文照原文排版。]
首先，考虑由 $z_1\in\bC$ 描述的运动自由度的贡献。超对称波函数是 $z_1$ 的全纯函数（乘以一个非全纯的高斯因子，该因子与我们的分析无关）。波函数为 $z_1^n$ 的态，在空间旋转下乘以 $e^{in\beta\epsilon_1}$。因此，在希尔伯特空间的这一部分上，该元素的迹为
\begin{equation}
\sum_{n=0}^\infty e^{in\beta\epsilon_1}=\frac1{1-e^{i\beta\epsilon_1}}.
\end{equation}
类似地，由 $z_2$ 描述的运动给出因子 $1/(1-e^{i\beta\epsilon_2})$。如果没有等同 $(u,v)\simeq(-u,-v)$，沿 $u,v$ 的运动也会类似地给出因子
\begin{equation}
\frac1{1-e^{i\beta((\epsilon_1+\epsilon_2)/2+a)}}\frac1{1-e^{i\beta((\epsilon_1+\epsilon_2)/2-a)}}.
\end{equation}
在非负整数 $n,m$ 对应的波函数 $u^n v^m$ 中，只有 $n+m$ 为偶数的波函数才与这一等同相容。因此，沿 $u$ 和 $v$ 运动的正确贡献为
\begin{equation}
\begin{aligned}
&\sum_{n\ge0,\ m\ge0,\ n+m:\text{偶数}}e^{in\beta((\epsilon_1+\epsilon_2)/2+a)}e^{im\beta((\epsilon_1+\epsilon_2)/2-a)}\\
&\qquad=\frac{1+e^{i\beta(\epsilon_1+\epsilon_2)}}{(1-e^{i\beta(\epsilon_1+\epsilon_2+2a)})(1-e^{i\beta(\epsilon_1+\epsilon_2-2a)})}.
\end{aligned}
\end{equation}
最后，不应忘记式 \eqref{constantshift} 中能量常数平移的贡献。

合起来得到单瞬子贡献
\begin{equation}
Z_1=e^{2\pi i\tau_{UV}}\frac1{1-e^{i\beta\epsilon_1}}\frac1{1-e^{i\beta\epsilon_2}}\frac{1+e^{i\beta(\epsilon_1+\epsilon_2)}}{(1-e^{i\beta(\epsilon_1+\epsilon_2+2a)})(1-e^{i\beta(\epsilon_1+\epsilon_2-2a)})}.\label{inst-z1}
\end{equation}

考虑如何取四维极限 $\beta\to0$。式 \eqref{inst-z1} 的分母为 $\beta^4$ 阶；因此，要使 $Z_1$ 有合理的极限，必须假设
\begin{equation}
e^{2\pi i\tau_{UV}}\sim\beta^4.
\end{equation}
令 $\beta^{-1}=\Lambda_{UV}$，可见这与 $\tau_{UV}$ 的单圈跑动 \eqref{lambdaSU2pure} 相容。确实，应当认为，在 $x_5$ 方向半径的倒数所对应的能标附近，五维描述变为有效的四维描述。我们决定在保持
\begin{equation}
\underline\Lambda^4:=\beta^{-4}e^{2\pi i\tau_{UV}}
\end{equation}
固定的同时取 $\beta\to0$ 极限。于是得到
\begin{equation}
Z_1\to\underline\Lambda^4\frac12\frac1{\epsilon_1}\frac1{\epsilon_2}\frac1{(\epsilon_1+\epsilon_2)/2+a}\frac1{(\epsilon_1+\epsilon_2)/2-a}.
\end{equation}
将其代入式 \eqref{main-instanton-formula}，得到预势到这一阶的瞬子贡献：
\begin{equation}
2\pi iF_\text{non-pert}=\lim_{\epsilon_{1,2}\to0}\epsilon_1\epsilon_2\log(1+e^{2\pi i\tau_{UV}}Z_1+\cdots)=\frac12\frac{\underline\Lambda^4}{a^2}+\cdots.
\end{equation}
我们再现了与式 \eqref{results-from-curve} 相同的系数 $d_1=1/2$，但还不能就此宣布成功，因为在这一粗略推导中，没有办法预先确定曲线一侧的瞬子因子 $\Lambda^4$ 与瞬子计算一侧的对应因子 $\underline\Lambda^4$ 相同。相反，这一结果是对 $\Lambda^4=\underline\Lambda^4$ 的确认。

\begin{figure}[h]
\[
\inc{hook}
\]
\caption{给定杨图中的方格 $s$ 时，臂长与腿长的定义。$A_Y(s)$ 和 $L_Y(s)$ 分别是白圆和黑圆的数目。\label{fig:armleg}}
\end{figure}
因此，要真正看到二者一致，就需要继续算下去。解释如何进行这一计算需要多得多的篇幅，但结果可以写成下面的组合公式：
\begin{multline}
\tilde Z_k=\sum_{Y_1,Y_2}\prod_{n,m=1}^2\prod_{s\in Y_n}\frac1{1-e^{i\beta(-L_{Y_m}(s)\epsilon_1+(A_{Y_n}(s)+1)\epsilon_2+a_m-a_n)}}\\
\prod_{t\in Y_m}\frac1{1-e^{i\beta((L_{Y_n}(t)+1)\epsilon_1-A_{Y_m}(s)\epsilon_2+a_m-a_n)}}.\label{inst-general}
\end{multline}
这里，求和遍历总方格数为 $k$ 的杨图对 $(Y_1,Y_2)$；$s\in Y$ 表示 $s$ 是杨图 $Y$ 中的一个方格；最后，函数 $A_Y(s)$ 和 $L_Y(s)$ 是杨图 $Y$ 中方格 $s$ 的臂长和腿长，其定义见图 \ref{fig:armleg}。

从这一一般公式再现式 \eqref{inst-z1} 给出的 $Z_1$，是一个很好的练习；有贡献的杨图对为 $(\Young{1},0)$ 和 $(0,\Young{1})$。计算 $Z_2$ 时，需要对五个杨图对 $(\Young{2},0)$、$(\Young{11},0)$、$(\Young{1},\Young{1})$、$(0,\Young{2})$ 和 $(0,\Young{11})$ 求和；你会发现，在 $\beta\to0$ 极限中，
\begin{equation}
Z_2\to\Lambda^8\frac{(8(\epsilon_1+\epsilon_2)^2+\epsilon_1\epsilon_2-8a^2)}{\epsilon_1^2\epsilon_2^2((\epsilon_1+\epsilon_2)^2-4a^2)((2\epsilon_1+\epsilon_2)^2-4a^2)((\epsilon_1+2\epsilon_2)^2-4a^2)}.
\end{equation}

再次将它代入式 \eqref{main-instanton-formula}，得到预势到这一阶的瞬子贡献：
\begin{align}
2\pi iF_\text{non-pert}&=\lim_{\epsilon_{1,2}\to0}\epsilon_1\epsilon_2\log(1+Z_1+Z_2+\cdots)\\
&=\lim_{\epsilon_{1,2}\to0}\epsilon_1\epsilon_2\left(Z_1+(Z_2-\frac12Z_1{}^2)+\cdots\right)\label{225}\\
&=\frac12\frac{\Lambda^4}{a^2}+\frac5{64}\frac{\Lambda^8}{a^6}+\cdots,\label{226}
\end{align}
其中，从式 \eqref{225} 到式 \eqref{226} 还取了 $\beta\to0$ 极限。确实再现了式 \eqref{results-from-curve} 中已经见到的 $d_2=5/64$。使用任意喜欢的计算机代数系统，都可以直接检验：由曲线确定的系数 $d_k$，即式 \eqref{results-from-curve}，与由瞬子公式 \eqref{inst-general} 确定的系数始终一致。例如，文献 \cite{Nekrasov:2003rj} 给出了一般的数学证明。
